\par\medskip\noindent\textbf{Cómo leer la salida.}
{\small\setlength{\tabcolsep}{4pt}\renewcommand{\arraystretch}{1.15}
\begin{xltabular}{\linewidth}{@{}>{\raggedright\arraybackslash}p{0.24\linewidth}>{\raggedleft\arraybackslash}p{0.15\linewidth}>{\raggedright\arraybackslash}X@{}}
\toprule Indicador & Valor & Interpretación \\ \midrule \endfirsthead
\toprule Indicador & Valor & Interpretación \\ \midrule \endhead
\bottomrule \endlastfoot
$T$ & 819 meses & Excesos de rendimiento mensuales de 1949 a 2017, en \%. \\
\texttt{media} & 0,344 a 1,155 & Exceso de rendimiento promedio mensual de cada cartera; crece de crecimiento (V1) a valor (V5) y es mayor en las pequeñas (S1). \\
CAPM: \texttt{alfa}, \texttt{t(alfa)} & $-0{,}547$ ($t=-3{,}17$) a 0,470 ($t=3{,}75$) & Rendimiento no explicado por el mercado; 6 de 9 son significativas al 5\,\%. \\
CAPM: \texttt{beta} & 0,85 a 1,38 & Las betas no explican el patrón: S1V1 tiene la mayor beta y el menor rendimiento. \\
CAPM: \texttt{R2} & 0,59 a 0,89 & El mercado explica la mayor parte de la variación, menos en las pequeñas. \\
FF3: \texttt{alfa} & $-0{,}533$ a 0,136 & Las alfas se reducen en la mayoría; la gran excepción es S1V1 (pequeñas de crecimiento), que sigue siendo muy negativa. \\
\texttt{b\_SMB} & $-0{,}25$ a 1,40 & Sensibilidad al factor tamaño: alta en las pequeñas, negativa en las grandes. \\
\texttt{b\_HML} & $-0{,}41$ a 0,84 & Sensibilidad al factor valor: crece de V1 a V5 en cada fila de tamaño. \\
FF3: \texttt{R2} & 0,82 a 0,95 & Los tres factores explican mucho más que el mercado solo. \\
Promedio de $\lvert$alfa$\rvert$ & 0,256 / 0,129 & El modelo de tres factores reduce el error de valoración medio a la mitad. \\
S1V5 $-$ S1V1 & 0,811 / 1,017 / 0,653 & Diferencia de medias y de alfas entre valor y crecimiento pequeñas: el CAPM no la explica (incluso la agranda); FF3 explica una parte, no toda. \\
\end{xltabular}}
